% MF45331C
% Hyperventilationssyndrom
% 14.0
% 2013-11-30

\label{m:Hyperventilationssysndrom} 
\EE{Taktik}:
\Speech{Beruhigen, Differentialdiagnosen ausschließen, symptomatische Therapie}

\begin{itemize}
  \item Psychische Betreuung: Patient beruhigen, mit ruhiger, langsamer Stimme
  mit dem Patienten sprechen
  \item Kommandoatmung
  \item \ce{CO2}-Rückatmung: In ein Plastiksackerl über kurze Zeit ein- und
  ausatmen lassen. Ebenso eignet sich eine \EE{\ce{O2}-Maske mit Reservoir} oder ein (puderfreier) 
  Handschuh.
  \item Grundsätzlich sollte kein \ce{O2}  verabreicht werden. Aus
  psychologischen Gründen kann es jedoch sinnvoll sein \EE{\ce{O2} in geringer
  Dosierung (1\LPerMin*{})} über eine \ce{O2}-Maske zu
  verabreichen.
  \item Wenn nicht erfolgreich und Krämpfe bestehen bleiben bzw.
  Bewusstseinsstörungen auftreten: \TxMassVitMK 
\end{itemize}